\chapter{Quasinormal Modes}
\label{chap:qnm}
\chapterepigraph{``Rest, rest, perturb\`ed spirit.''}{William Shakespeare, \textit{Hamlet}, Act I, Scene 5}

\section{Three equivalent definitions}

The Green kernel obtained in the previous chapter is
\begin{equation}
G(x,x';\omega)
=
-\frac{f_-(x_<,\omega)f_+(x_>,\omega)}{\vW(\omega)}.
\label{eq:qnm-start-green}
\end{equation}
When $\vW(\omega)$ vanishes, the solution satisfying the outgoing condition on the left and the solution satisfying the outgoing condition on the right become the same solution.

\begin{definition}[Quasinormal mode]
A complex frequency $\omega_n$ satisfying
\begin{equation}
\vW(\omega_n)=0
\label{eq:qnm_wronskian_zero}
\end{equation}
is called a \term{QNM frequency}. The corresponding nontrivial homogeneous solution $\psi_n(x)$ satisfies radiation conditions directed out of the system at both asymptotic ends.
\end{definition}

With the time factor $\rme^{-\rmi\omega t}$, the conditions in the black-hole exterior are
\begin{align}
\psi_n(x)
&\sim
\rme^{-\rmi\omega_n x},
&&x\longrightarrow-\infty,
\label{eq:qnm_bc_left}
\\
\psi_n(x)
&\sim
\rme^{+\rmi\omega_n x},
&&x\longrightarrow+\infty.
\label{eq:qnm_bc_right}
\end{align}
Outgoing through the left endpoint means entering the event horizon. Thus a \term{QNM} has three equivalent characterizations: a solution of the \term{radiation boundary-value problem}, a zero of the Jost Wronskian, and a pole of the analytically continued Green function.

\section{Simple poles and residues}

If $\vW'(\omega_n)\ne0$, then $\omega_n$ is a \term{simple pole}, and
\begin{equation}
\mathop{\mathrm{Res}}_{\omega=\omega_n}
G(x,x';\omega)
=
-\frac{
f_-(x_<,\omega_n)f_+(x_>,\omega_n)
}{
\vW'(\omega_n)
}.
\label{eq:qnm_green_residue}
\end{equation}
Multiplying a Jost solution by a constant rescales numerator and denominator by the same factor, so this \term{residue} is independent of normalization.

The \term{QNM} waveform $\psi_n$ itself, however, can be multiplied by an arbitrary constant. Let $\bra{O}$ be the observation operation, $\ket{S}$ the source, and $A_n$ the \term{operator-valued residue}. The observed pole contribution is
\begin{equation}
\vA_n(\omega)
=
\frac{a_n}{\omega-\omega_n},
\qquad
a_n
=
\bra{O}A_n\ket{S}.
\label{eq:qnm-observable-residue}
\end{equation}
The background-determined frequency, a normalization-dependent waveform, the operator-valued \term{residue}, and an amplitude specified by a source and observer are not the same quantity \cite{Motohashi:2024fwt}.

\begin{principlebox}
Frequencies alone do not determine a time-domain waveform. The physical \term{mode strength} is the \term{matrix element} obtained by sandwiching the Green-function residue between the source and the observation operation.
\end{principlebox}

\section{Spatial divergence and temporal decay}

For a damped QNM of a stable black hole, $\im\omega_n<0$. At a fixed position,
\begin{equation}
\left|\rme^{-\rmi\omega_n t}\right|
=
\rme^{(\im\omega_n)t}
\end{equation}
decays. But at the right endpoint the spatial factor satisfies
\begin{equation}
\left|\rme^{+\rmi\omega_n x}\right|
=
\rme^{-(\im\omega_n)x},
\end{equation}
which grows as $x\to+\infty$; the same happens at the left endpoint. This is the result of extending an outward-propagating wavefront over all space with a single complex frequency. It is not a temporal instability.

\begin{warningbox}
A QNM is not an ordinary $L^2$ eigenfunction. If its divergence on the real axis is ignored and the QNMs are treated as a \term{complete set} in a \term{Hilbert space}, both normalization and the domain of a mode expansion can be misidentified. The \term{complex scaling} defined in Chapter~\ref{chap:complex-scaling} converts this divergence into decay along an analytically continued coordinate.
\end{warningbox}

\section{Energy fixes the direction of damping}

For a real potential $V(x)\geq0$ and a real solution of \term{finite energy}, define the energy on $[-R,R]$ by
\begin{equation}
E_R(t)
=
\frac{1}{2}
\int_{-R}^{R}\rmd x
\left[
(\partial_t\Psi)^2
+(\partial_x\Psi)^2
+V\Psi^2
\right].
\label{eq:qnm-energy-estimate}
\end{equation}
Using the wave equation and integrating by parts gives
\begin{equation}
\frac{\rmd E_R}{\rmd t}
=
\left[\partial_t\Psi\,\partial_x\Psi\right]_{-R}^{R}.
\label{eq:qnm-energy-flux-balance}
\end{equation}
If only right-moving waves are allowed at the right endpoint and left-moving waves at the left endpoint, then as $R\to\infty$ the right-hand side is the negative flux escaping through the two ends, and the exterior energy cannot increase.

Suppose there were an outgoing mode with $\im\omega>0$. Its spatial profile would decay at both ends and have finite energy, while its time factor would grow exponentially. This contradicts Eq.~\eqref{eq:qnm-energy-flux-balance}. Hence this positive-energy problem has no \term{unstable QNMs}. The argument does not determine the damping rate and cannot be applied unchanged when $V$ is not positive, when rotation permits negative horizon flux, or when gravitational constraints must be included. Stability is not an empirical pattern in a frequency table; it is a question of proving the sign of an appropriate energy and its \term{boundary flux}.

\section{A solvable barrier}

To see the origin of QNMs in one formula, consider the \term{P\"oschl--Teller barrier}
\begin{equation}
V(x)
=
\frac{V_0}{\cosh^2(\alpha x)},
\qquad
V_0>0,
\quad
\alpha>0.
\label{eq:poschl_teller_potential}
\end{equation}
With the change of variable $y=(1+\tanh\alpha x)/2$, the equation becomes a \term{hypergeometric equation}. Imposing the radiation conditions at both ends gives
\begin{equation}
\omega_n
=
\mathord{\pm}
\sqrt{V_0-\frac{\alpha^2}{4}}
-\rmi\alpha\left(n+\frac{1}{2}\right),
\qquad
n=0,1,2,\ldots.
\label{eq:poschl_teller_qnm}
\end{equation}
The barrier height controls the real part, while its width controls the spacing of damping rates. This is not the exact black-hole potential, but it shows that open boundary conditions alone generate complex frequencies.

\section{The photon-sphere limit}

For Schwarzschild spacetime with $\ell\gg1$, the leading part of the effective potential is $f(r)\ell(\ell+1)/r^2$, whose maximum lies at the \term{unstable photon orbit} $r=3M$. The angular frequency and \term{Lyapunov exponent} of that orbit are both
\begin{equation}
\Omega_c
=
\Lambda_c
=
\frac{1}{3\sqrt{3}M}.
\end{equation}
For fixed \term{overtone} number $n$,
\begin{equation}
\omega_{\ell n}
\simeq
\Omega_c\left(\ell+\frac{1}{2}\right)
-\rmi\Lambda_c\left(n+\frac{1}{2}\right).
\label{eq:eikonal_qnm}
\end{equation}
Local ray dynamics explains the slowly damped part of the spectrum, but finite-$\ell$ poles, residues, and the \term{branch cut} are not obtained from this approximation. For example, the fundamental $\ell=2$ gravitational mode is
\begin{equation}
M\omega_{20}
\simeq
0.37367-0.08896\rmi.
\label{eq:schwarzschild_fundamental_qnm}
\end{equation}
\cite{Berti:2009kk}

\section{Diagnostics for numerical calculations}

Direct integration finds zeros of the Wronskian of the left and right Jost solutions. The \term{continued-fraction method} imposes a \term{minimal-solution condition} on an \term{asymptotic series}. Time evolution estimates poles from the waveform. The complex-scaling method introduced in Chapter~\ref{chap:complex-scaling} converts outgoing waves into decaying boundary conditions. Although the methods differ, their diagnostics are common.

\begin{enumerate}[label=(\roman*),leftmargin=2.8em]
\item The time convention and the radiation conditions at both ends must be consistent.
\item The poles must remain stable when the basis size, domain, and numerical precision are varied.
\item Discretized continuous components must be separated from genuine poles.
\item If a physical amplitude is required, the residue must converge as well.
\end{enumerate}

Complex-scaling calculations for Schwarzschild and Reissner--Nordstr\"om are given in Ref.~\cite{Ogawa:2026veu}. In a \term{non-self-adjoint problem}, sensitivity of eigenvalues must be distinguished from sensitivity of the real-frequency response \cite{Jaramillo:2020tuu}.

\section{Conclusion of this chapter}

A QNM is not a material vibration intrinsic to a black hole. It is a resonance pole of a Green function with open boundary conditions. The pole is important, but a response also requires the residue, source, observation operation, and continuous components. In the next chapter we invert the pole and continuous contributions and ask which parts of the time-domain waveform they control.

\begin{researchproblem}[Numerical implementation of Kerr QNMs]
Solve the angular and radial Teukolsky equations simultaneously and track the fundamental QNM with $s=-2$ and $\ell=m=2$ from $a/M=0$ into the near-extremal regime \cite{Teukolsky:1972my,Leaver:1985ax}. At minimum, construct an implementation satisfying the following requirements.

\begin{enumerate}[label=(\roman*),leftmargin=2.8em]
\item Determine the angular eigenvalue and $\omega$ in the same iteration or \term{contour method}.
\item Vary the series order, numerical domain, and precision to separate stable poles from discretized continuous components.
\item Verify that the Schwarzschild value is recovered as $a\to0$.
\item If possible, compute the Green-function residue and identify quantities invariant under a change of normalization.
\end{enumerate}

Reproducing a frequency table is only the entry point. Formulate, from the implementation itself, a research problem concerning at least one of the following: the near-extremal \term{branch structure}, the \term{zero-damped mode}, or the source-dependent \term{excitation amplitude}.
\end{researchproblem}
